\section{Étude et ré-implémentation simplifiée du modèle \gls{soa}}
\subsection{Travail effectué}\label{subsec:codebase}
La deuxième partie du projet à été la prise en main de la base de code fournie.
Il s'agit d'une implémentation \gls{soa} d'un \gls{lm} au niveau du caractère, sur laquelle notre maître de stage avait commencé à travailler.
Le code d'origine provient du dépôt \og awd-lstm-lm\fg{}\autocite{awd_source}, qui contenait un modèle \gls{soa} de \gls{lm} au niveau du caractère.

Au début du stage, la base de code contenait~:
\begin{itemize}
	\item la version d'origine du dépôt~;
	\item un début de ré-implémentation simplifiée du \gls{model} de la version d'origine~; cette version  devait servir de base pour développer le \gls{model} du \gls{gmsnn}, ainsi que de comparaison pour les performances du nouveau \gls{model}~; elle comportait quelques \glspl{bug} et ne fonctionnait pas en l'état~;
	\item un début de travail sur l'architecture du \gls{gmsnn}.
\end{itemize}

\vspace{1em}
L'objectif de cette étape était de faire fonctionner la ré-implémentation simplifiée du \gls{model}.

Pour cela, nous avons déchiffré et re-documenté le code, qui contenait des fragments obsolètes et peu documentés.
Après le déchiffrage, il à fallu comprendre et réparer les fragments défectueux.

\subsection[\Glsentrytext{model} ré-implémenté simplifiée]{\Gls{model} ré-implémenté simplifiée}
Le \gls{model} simplifiée produit est composé d'un module encodant les caractères, d'un \gls{rnn} particulier (un \gls{lstm}, voir \autoref{def:lstm}), et d'un module produisant une distribution de probabilité sur les caractères connus. La \autoref{fig:reimplement} représente cette architecture.

\begin{figure}[h]
	\centering
	\scalebox{1}{\usetikzlibrary{calc}
\newsavebox\mydictbox
\savebox\mydictbox{%
  \includegraphics[width=3em]{/home/emarquer/InternshipSynalp2018/report_internship/images/tool_translate_offline_dictionary.png}
}

\begin{tikzpicture}[shorten >=2pt,->,draw=black, node distance=7em]
    \tikzstyle{every pin edge}=[<-,shorten <=2pt]
	\tikzstyle{module}=[minimum size=3em, fill=gray!50!]
	\tikzstyle{char}=[module, circle, fill=green!50!]%, label={below:\usebox\mydictbox}]
	\tikzstyle{text label}=[rectangle, text centered, text width=7em, node distance=5em]
	\tikzstyle{nn}=[rectangle, module, fill=orange!50!]
	\tikzstyle{rnn}=[nn, fill=red!50!]

	\node[char] (char) {\$};
	\node[nn, right of=char] (emb) {Emb.};
	\node[rnn, right of=emb] (rnn) {LSTM};
	\draw[->] (rnn.north) to [out=90,in=90, looseness=2] ($(rnn.west) + (-0.5,0)$) to [out=-90,in=-90, looseness=2] (rnn.south);
	\node[nn, right of=rnn] (lin) {Lin.};
	\node[char, right of=lin] (out) {P(\$$_{+1}$)};
	\draw[->] (char) to (emb);
	\draw[->] (emb) to (rnn);
	\draw[->] (rnn) to (lin);
	\draw[->] (lin) to (out);
	%\node[char] (char) {\$};


    \node[text label, above of=char] (charl) {Caract\`{e}re};
    \node[text label, above of=emb] (embl) {Encodage du caract\`{e}re};
    \node[text label, above of=rnn] (rnnl) {RNN};
    \node[text label, above of=lin] (linl) {Transformation en distribution de probabilit\'{e}};
    \node[text label, above of=out] (outl) {Probabilit\'{e} d'apparition des caract\`{e}res};
\end{tikzpicture}}
	\caption[Architecture du \glsentrytext{model} réimplémenté]{Architecture du \gls{model} réimplémenté. Le modèle prend en entrée des caractères, et produit des probabilités sur quel caractère apparaîtra ensuite.}\label{fig:reimplement}
\end{figure}

Le module d'encodage des caractères, appelé \foreign{embeding layer} en anglais (littéralement \og couche d'inclusion\fg{}), produit une représentation apprise de chaque caractère sous forme de \gls{tensor}. Ce \gls{tensor} est appelé \gls{embedding}. Ce module, entraîné, peut apprendre des propriétés spécifiques à chaque caractère. Par exemple ce module peut apprendre que telle lettre est une consonne et que telle autre est un caractère de ponctuation. \label{def:embeding}

Le \gls{rnn} traite les caractères sous forme de séquence, et peut ainsi apprendre la structure en mots, la syntaxe et d'autres propriétés du langage.

Le module produisant la distribution de probabilité est un module linéaire\footnote{\og Module linéaire \fg{} est le nom donné aux modules composé d'un \gls{nn} intégralement connecté. Un réseau de neurones intégralement connécté signifie que chaque neurone d'une couche est connecté avec tous les neurones de la couche précédente. Il s'agit de l'architecture de \gls{nn} la plus simple.\label{def:fully_connected}\label{def:lin_module}}.
Il transforme les informations produites par le réseau de neurones en probabilité pour chaque caractère d'être le prochain caractère de la séquence.

Voir les annexes \ref{subsec:testreimp_1}, \ref{subsec:testreimp_2} et \ref{subsec:testreimp_3} pour les rapports sur le \gls{model} ré-implémenté. 

\subsection{Conclusion}
Cette étape, outre la préparation du code, à permit la mise en œuvre et une meilleur compréhension des concepts appris durant l’étape précédente. C'est la première pierre de l'édifice qu'est le \gls{gmsnn}.